\section{Four boundary labels and convex assembly}\label{sm:masks}

An open cell has one ordered quadrilateral gluing face
$Q=(C_0,C_1,C_2,C_3)$. Number its edges $j=(j,j+1)$ modulo four.
Edges $0,2$ are called free, and $1,3$ cap edges. These names record
the original leaf order, not an unmarked symmetry of the whole tangle.
Set $\eta_j=1$ if the face beside edge $j$, other than $Q$, is shaded,
and set $\eta_j=0$ if it is white. The four bits
$\boldsymbol\eta=(\eta_0,\eta_1,\eta_2,\eta_3)$ will be called the
boundary mask.

\begin{proposition}\label{sm:mask-invariant}
At every open subexpression evaluated at its strict stationary wires
there is a convex ideal cell with a maximal, nondegenerate ordered
quadrilateral output. Its boundary-edge angles are
\begin{equation}\label{sm:masked-angles}
 \frac{\theta_j}{\pi}=
 \begin{cases}
 b-\eta_j/2,&j=0,2,\\
 1-b-\eta_j/2,&j=1,3.
 \end{cases}
\end{equation}
Every other surviving edge is right angled and separates a white
face from a shaded triangle. The possible supports satisfy
\begin{equation}\label{sm:mask-support}
 L:\eta_0=\eta_2=0,\qquad H:\eta_1=\eta_3=0,\qquad
 A:\boldsymbol\eta=0.
\end{equation}
Every original shaded triangle remains an entire face. The four
output corners are four active original terminal singletons.
Every other main ideal vertex is a permanently consumed pair of
terminals; crossing-circle apices stay separate.
\end{proposition}

The statement permits a type $H$ output with only one shaded free
edge. Replacing its two free angles by a common angle would therefore
lose information essential to the next gluing.

\begin{figure}[ht]
\centering\begin{tikzpicture}[x=1cm,y=1cm,line cap=round,
  every node/.style={font=\small},>=Stealth]
 \definecolor{shadeedge}{RGB}{153,61,37}
 \foreach \x/\title/\bits in {0/$F:\ H$/{},4.7/$G:\ L$/{},9.4/$FG:\ H$/{}} {
  \begin{scope}[shift={(\x,0)}]
   \coordinate (a) at (0,0);
   \coordinate (b) at (2,0);
   \coordinate (c) at (2,1.6);
   \coordinate (d) at (0,1.6);
   \draw[gray!65] (a)--(b)--(c)--(d)--cycle;
   \foreach \p/\lab/\pos in {a/0/below left,b/1/below right,c/2/above right,d/3/above left}
    {\fill (\p) circle (1.5pt); \node[\pos] at (\p) {$C_{\lab}$};}
   \node at (1,2.25) {\title};
  \end{scope}
 }
 \draw[shadeedge,very thick] (0,0)--(2,0) (0,1.6)--(2,1.6);
 \draw[shadeedge,very thick] (4.7,0)--(4.7,1.6) (6.7,0)--(6.7,1.6);
 \draw[shadeedge,very thick] (9.4,0)--(11.4,0);
 \node at (1,-.75) {$(1,0,1,0)$};
 \node at (5.7,-.75) {$(0,1,0,1)$};
 \node at (10.4,-.75) {$(1,0,0,0)$};
 \node at (3.35,.8) {$\cdot$};
 \draw[->,thick] (7.5,.8)--(8.6,.8);
 \node[anchor=west,font=\footnotesize] at (0,-1.5)
 {$\eta^o_0=\eta^F_0,\quad \eta^o_2=\eta^G_2,\quad
   \eta^o_1=\eta^o_3=0$};
\end{tikzpicture}

\caption{Individual boundary labels in an $HL$ product. A colored
edge has a shaded adjacent face; the other edges have a white
adjacent face. The turned leaf has mask $(1,0,1,0)$, the upper
leaf has mask $(0,1,0,1)$, and the product has mask $(1,0,0,0)$.
Its two free angles are $\pi(z-1/2)$ and $\pi z$.
The pictures display the ordered quadrilateral gluing faces.}
\label{sm:mask-figure}
\end{figure}

\subsection{Every seam angle}
For a leaf the mask is $(0,1,0,1)$, so
\eqref{sm:masked-angles} is the pyramid angle calculation.
A turn rotates the mask by one and replaces $b$ by $1-b$.
This proves the assertion under turns and exchanges the support
conditions for $L,H$.

At a product the output mask is exactly
\begin{equation}\label{sm:mask-product}
         \boldsymbol\eta^o=(\eta^F_0,0,\eta^G_2,0).
\end{equation}
Indeed the two output cap sides adjoin the prism's white triangles,
and each output free side inherits its indicated child side.
A child cap seam has angle
$\pi(1-\eta_j/2)$: a zero bit merges two white pieces flatly, while
a one bit leaves a right shaded--white edge.
At the edge common to both prism inputs, the total angle is
\begin{equation}\label{sm:shared-angle}
 \pi\left[w+x-\frac{\eta^F_2}{2}
             +y-\frac{\eta^G_0}{2}\right]
       =\pi\left[1-\frac{\eta^F_2+\eta^G_0}{2}\right].
\end{equation}
Both bits could be one only with two type $H$ inputs, forbidden
by \eqref{sm:types}. Thus this seam is either flat white--white
or right shaded--white. The remaining free output angles are
$\pi(p+x-\eta^F_0/2)$ and $\pi(q+y-\eta^G_2/2)$; the cap
angles are $\pi(1-z)$. This gives \eqref{sm:masked-angles}.

The support assertion also follows case by case from
\eqref{sm:types}. An output of type $A$ has no type $H$ child,
so both inherited free bits vanish. An output of type $H$ has
zero cap bits by \eqref{sm:mask-product}. Turns give all type $L$
cases. We assert these support inclusions, not equality with the
two-bit leaf masks.

For successive convex attachments one must check the intermediate
stage, not only the final totals. Attach $F$ to the prism first.
Its free seams have angles
\[
 \pi(z-\eta^F_0/2),\qquad
 \pi(1-y-\eta^F_2/2).
\]
If the first bit is one, $F$ has type $H$, hence $z>x>1/2$.
If the second bit is one, $x>1/2$ and $x+y<1$ give $y<1/2$.
Both displayed angles are therefore positive and strictly less
than $\pi$. In particular the other prism input remains a
supported nondegenerate quadrilateral. The cap seams have already
been checked. The second attachment is now justified by
\eqref{sm:shared-angle}.

\subsection{Convexity includes the ideal boundary}
Here is the geometric argument used at each attachment. Place two
convex finite-faced ideal polyhedra on opposite sides of the plane
of the matched quadrilateral. Their finite intersection is that
face. The finite union is closed, connected and locally convex:
the only new boundary points lie on the seam edges, where the
total angle is at most $\pi$. There are no finite seam vertices.

Two finite points in the union have a shortest path in it.
Indeed a length-bounded minimizing sequence stays in a compact
ambient hyperbolic ball. Constant-speed parametrization and
Arzelà--Ascoli give a limiting path; closedness keeps it in the
union. Local convexity makes each small subarc the ambient
geodesic. The complete path is therefore the unique ambient
geodesic segment. This proves convexity of the union.

There is a further ideal-point issue. The two ideal closures can
meet only on the boundary circle of the common supporting plane.
The intersection of either cell with that plane is exactly its
maximal quadrilateral face. An additional ideal vertex on the
circle would enlarge that face past one of its edges, making the
edge flat; all four supporting angles there are strictly nonflat.
Equivalently, in the Klein model the closed quadrilateral meets
the sphere at infinity only in its four vertices. Thus the only
ideal coincidences are the four specified corner identifications.
This argument applies again to the second input, which the
intermediate angle check kept strictly supported.

Every output angle in \eqref{sm:masked-angles} lies in $(0,\pi)$
on its actual type interval. Hence the output remains a maximal
supported quadrilateral and the induction continues. Every
shaded-triangle seam is right, never flat. Only white faces merge:
no shaded edge or triangle is deleted, split or merged with
another shaded face. A leaf apex is never in an attaching
quadrilateral. It retains its four right edges and its two
original shaded sectors.

\subsection{Terminal pairs never become chains}
Give leaf $i$ the four distinct terminal symbols $(i,0),\ldots,(i,3)$
and the distinct apex symbol $O_i$. The purely literal endpoint
rules prove the following invariant independently of the metric.
An open expression has four active singleton terminals; all other
terminals form disjoint consumed two-element pairs. A turn only
permutes the active tuple. Product consumes exactly
\begin{equation}\label{sm:terminal-product}
          F_2\sim G_1,\qquad F_3\sim G_0,
\end{equation}
and keeps $F_0,F_1,G_2,G_3$ active. The child leaf sets are disjoint.
A consumed pair is never subsequently attached or exposed.
Thus no class grows to three terminals. The geometric maps
\eqref{sm:stack-map} make exactly these pairs; the ideal-point
argument excludes any additional identification.

At an exposed corner the two adjacent mask bits have sum at most one.
If their sum is zero, its cyclic faces are the gluing quadrilateral,
white, shaded, white; if the sum is one there are only the
quadrilateral, white and shaded faces in the retained order.
This also follows by taking the ideal Euclidean link:
the two quadrilateral-edge angles sum to $\pi$ or $\pi/2$,
and the remaining edge angles are right. Each exposed corner
therefore carries exactly one original shaded sector.
At \eqref{sm:terminal-product}, the two sectors concatenate
in the original ribbon order. This proves every assertion of
Proposition~\ref{sm:mask-invariant}.

\section{The closed polyhedron and the original link}\label{sm:cut}

At the closing root glue the entire face by
\begin{equation}\label{sm:closing}
                       C_j^F\longmapsto C_{3-j}^G.
\end{equation}
The profile equation already matches its ordered metric.
An edge $j$ on the first child matches edge $2-j$ on the second.
Its total angle is $\pi$ minus half the sum of the two mask bits,
in units of $\pi$. Two shaded free bits would require an $HH$
closing pair; two shaded cap bits would require $LL$. Neither is
strict. The convex-union proof therefore applies, and all remaining
edges of the closed polyhedron $P_E$ are right shaded--white edges.

The four closing pairs in \eqref{sm:closing} are between distinct
root-child leaf sets. Together with \eqref{sm:terminal-product}
they pair every original terminal once, never with a terminal of
the same leaf. Hence each leaf still has four distinct corner
vertices. Its two shaded triangles are distinct and meet only
at their original apex. This rules out a transitive corner
collapse at closure.

The boundary of the actual convex polyhedron is a sphere. Its
white faces are maximal convex ideal polygons. They cannot be
bigons or zero-width regions. No two shaded triangles share an
edge: every surviving edge has a white side. These observations
will matter for both the topological cut and the cusp returns.

\subsection{Recovering every colored face from the literal wires}
Compile the same boundary without prisms. At leaf $i$ put the
two shaded triangles
\[
       (O_i,C_{i1},C_{i2}),\qquad (O_i,C_{i3},C_{i0})
\]
in its original oriented projection chart. At a terminal record
its outgoing-next and incoming-previous shaded neighbors. Joining
two actual identity-wire endpoints concatenates those two ordered
sectors. To see the order directly, take a thin planar neighborhood
of the joining arc: its two white sides join the complementary
sectors, with no crossing or reflected child.
Turns rotate the charts; the final closure uses
\eqref{sm:closing}. This constructs the full oriented ribbon graph,
not merely its abstract incidence or component count.

All white face cycles are now determined. A facial walk entering
along $(u,v)$ leaves $v$ along the neighbor preceding $u$ in its
cyclic order. The original shaded triples are some of these walks;
the remaining walks are exactly the white regions.
The prism maps \eqref{sm:stack-map} induce the same two pairings
and the same sector concatenation. The mask proof shows that
every other seam is flat white--white and every shaded edge remains.
Thus the metric boundary has precisely this ribbon graph and
every one of its colored facial cycles.

\subsection{The reversible disk cut}
Each original auxiliary circle lies in its own leaf ball, perpendicular
to the projection sphere, with its recorded twice-punctured disk.
The disks and their collars remain disjoint under physical turns
and products. In either projection ball each auxiliary semicircle
is boundary parallel by the corresponding half disk.
Drilling the semicircle introduces one handle; cutting along that
same bridge half disk cancels the handle in its explicit collar.
The disjoint cancellations leave a ball. The main segments lie
on the projection boundary; removing their neighborhoods creates
ideal ends and no interior handle.

The cut ball has one ideal vertex for each paired main segment
and one for each auxiliary semicircle. The exposed half disks give
exactly the two shaded triangles above. The remaining regions
have precisely the compiled white facial walks. Identify this
labeled boundary sphere with $\partial P_E$ and extend over the
balls by radial extension after choosing interior points.
Across the two balls restore the identically labeled white
projection faces. For a pure leaf restore the shaded pairing
fixing $O_i$ and sending $C_{i1}$ to $C_{i0}$ and $C_{i2}$ to
$C_{i3}$. For an odd leaf use the original half-disk rotation,
which pairs the corresponding triangles across copies.
Reversing these cuts gives exactly $Y_E^\eps$.

This is a degree-one topological identification with all components
and paths retained. It does not interpret an intermediate gluing
quadrilateral as the physical four-punctured boundary of a tangle:
that assertion is unnecessary and generally false.

\subsection{Complete returns and a common horosphere decoration}
Take two oppositely oriented copies of $P_E$. For a pure leaf, an
edge $e$ of one shaded triangle has a distinct partner $e'$ on
the other. Its four occurrences are
\[
                e^{(0)},\ e'^{(0)},\ e'^{(1)},\ e^{(1)}.
\]
The labeled shaded maps are inverse maps and the white maps identify
equal labels across copies, so the return is pointwise identity
on the edge. The angle sum is $4(\pi/2)=2\pi$.
An odd pairing changes the copy permutation in this return.
It still has four distinct occurrences, since the two shaded
triangles have no common edge and the two cut copies are distinct.
The local pure and odd return identities of
\cite[equations~(5.7)--(5.8)]{Geometry}, written again in
\eqref{sm:edge-words} below, show that the full return is identity.
Thus the quotient has neither an edge cone angle nor a residual
translation along an edge.

At an ideal vertex all edges are right. Its small horosphere
section is a finite Euclidean polygon with all angles $\pi/2$,
hence a rectangle. Its side lengths are positive: the supporting
planes are distinct, and coincident opposite sides would force
the three-dimensional polyhedron into a plane. The colors alternate.
Scale its horosphere until its two equal shaded side lengths are one.
Each shaded ideal triangle now has its canonical three unit
horocyclic segments, so every shaded isometry preserves these
decorations, including the odd permutations. White maps preserve
them because they identify the identical labeled copies.
The rectangles therefore glue without dilation into smooth compact
Euclidean tori, the original cusp tori of the cut.
A common shrinking factor makes the horospheres disjoint; beyond
them are complete horoball ends. The truncated core is compact.
The resulting finite-volume metric is complete.

All cell attachments had disjoint interiors. There is one leaf
pyramid and one prism per retained gate in each cut copy. Their
volumes are $\pi g$ and $\pi m$, respectively. Consequently
$\Vol(Y_E^\eps)=2\Vol(P_E)=2\pi F_0$ as an equality of ordinary
metric volumes. This proves the metric part of
Theorem~\ref{sm:complete}, including the two-leaf strict case.

\section{The primitive collar and its component information}
\label{sm:collar}

The auxiliary apex retains its original leaf rectangle throughout
the construction. In \eqref{sm:leaf} its adjacent sides have lengths
$2\sin\pi a$ and $2\cos\pi a$. The cusp has exactly two such
rectangles. The original disk paths in
\cite[Section~7, equations~(7.1)--(7.5)]{Geometry} are local to this
unchanged collar. We recall the path and area checks so that no
global completed-root marking is being imported.

Normalize the ordered leaf base to $(0,\infty,1,R)$.
Let $W_j$ be the white transition at edge $j$ between projection
charts, and let $S,A,B$ be the pure and the two odd shaded maps,
as defined in Section~\ref{sm:maps}. The traced paths are
\begin{equation}\label{sm:peripheral-words}
 \lambda=W_2^{-1}W_0,\qquad
 \mu_0=S^{-1},\qquad
 \mu_-=A^{-1}W_0,\qquad
 \mu_+=W_0^{-1}B.
\end{equation}
The longitude crosses white face $0$ and returns through $2$.
The pure meridian crosses the shaded identification once in the
reverse direction. The odd paths cross the respective half-disk
pairing and white face $0$. The retained negative crossing uses
$\mu_-$ and the positive crossing uses $\mu_+$ in the original
disk frame. A physical turn transports the whole collar, not its sign.

Conjugation back to the original apex at infinity gives the translations
\begin{equation}\label{sm:translations}
 \mu_0=E-1,\qquad \lambda=-2(E+1),\qquad
 \mu_-=-2,\qquad \mu_+=2E.
\end{equation}
The actual torus area is $8\sin\pi a\cos\pi a=4\sin2\pi a$,
and
$|\Ima(\overline{\mu_0}\lambda)|=4\sin2\pi a$.
Thus the closed paths span the full lattice, not a proper sublattice.
Replacing $\mu_0$ by either odd meridian preserves this area.
They too are primitive. As translation vectors on the corresponding
half-twist partners,
\[
                   \mu_\eps=\mu_0-\eps\lambda/2.
\]
This compares different marked tori; half a longitude is not being
declared an integral curve on the pure torus.
Since $(E+1)/(E-1)=-\ii\cot\pi a$, the ratio
$\lambda/\mu_\eps$ is \eqref{sm:tau}.
This proves the remaining marking assertion of
Theorem~\ref{sm:complete}.

\begin{proposition}\label{sm:knot-choice}
Every strict expression admits a choice of pure or odd leaves for
which the main strands form one component. Either odd sign may be
chosen, and subsequent full twists preserve the component count.
\end{proposition}
\begin{proof}
Write the endpoint pairings as
$I=(03)(12)$, $U=(01)(23)$ and $X=(02)(13)$.
A leaf offers $\{U,X\}$, and a turn exchanges $I,U$ and fixes $X$.
Tracing the two actual product connections gives
\[
\begin{array}{c|ccc}
 &I&U&X\\\hline
 I&I&U&X\\
 U&U&U\text{ and an internal circle}&U\\
 X&X&U&I
\end{array}.
\]
Discard only the product $U\cdot U$ in this existence argument.
For sets $A=\{I,U\}$, $B=\{I,X\}$ and $C=\{U,X\}$, the output
sets for $AA,AB,AC,BB,BC,CC$ are respectively
$A,\{I,U,X\},C,B,\{I,U,X\},A$, also in reversed orders.
Thus every subexpression has at least two loop-free choices,
starting with the leaf itself; the same is true if a child already
offers all three. At least one choice at the root stack is $U$
or $X$, both of which close to one component. Backtracking chooses
the leaves. No identity of unmarked tangles or assertion about
quantum connected-sum factors is used.
\end{proof}
